\chapter{\texorpdfstring{Derived Polynomial Transposes of Contextual Mealy Execution and Guarded Program Structure}{Derived Polynomial Transposes of Contextual Mealy Execution and Guarded Program Structure}}
\label{ch:derived-polynomial-contextual-mealy}

\section{Orientation: polynomial transposes of cartesian execution}
\label{sec:derived-poly-orientation}

The previous derived chapters developed the homotopy-invariant Mealy calculus
inside the localized homotopy theory
\[
\mathcal C
=
\mathbf{Cat}(\mathcal E)_{JT,\infty}
=
L_{DK}(\mathbf{Cat}(\mathcal E),W_{JT})
\]
of internal categories.  A derived Mealy morphism from
\[
\mathbb A
\]
to
\[
\mathbb B
\]
was defined as a derived correspondence
\[
\chi:\mathbb X\to\mathbb A\times\mathbb B
\]
in
\[
\mathcal C.
\]
Writing its two components as
\[
\pi:\mathbb X\to\mathbb A,
\qquad
\Omega:\mathbb X\to\mathbb B,
\]
the map
\[
\Omega:\mathbb X\to\mathbb B
\]
is the derived \(\mathbb B\)-valued output \(1\)-cocycle.

The derived relative chapter showed that a downstream behaviour
\[
\rho:\mathbb Y\to\mathbb B
\]
contextualizes \(\chi\) by the homotopy pullback
\[
\mathbb Z
:=
\mathsf{Exec}^h_{\chi,\rho}
=
\mathbb X\times_{\mathbb B}\mathbb Y.
\]
Dynamic proarrow logic is then the logic of endocorrespondences
\[
\mathbb Z\nrightarrow\mathbb Z.
\]

The present chapter explains the polynomial form of this same execution
semantics in the cartesian part of the derived Mealy calculus.  The polynomial
transpose is not attached to an arbitrary derived correspondence.  It is
attached to right-fibrational derived Mealy data, where the input projection
carries homotopy-coherent contravariant transport along input arrows.

Thus the guiding principle is:
\[
\boxed{
\text{derived polynomial semantics is the dependent-product transpose of
cartesian contextual Mealy execution.}
}
\]

Equivalently:
\[
\boxed{
\text{proarrow semantics is flattened execution semantics, while polynomial
semantics is curried total-response semantics.}
}
\]

The strict polynomial chapter is recovered by choosing split internal
discrete-fibration presentations and then passing to the objectwise shadow of
the resulting homotopy-coherent arrow and response objects.  In such split
presentations, derived right fibrations become ordinary right actions, derived
cartesian transport becomes the strict action map, and the derived response
coalgebra has the familiar objectwise formula
\[
(x,y)
\longmapsto
\left(
a
\longmapsto
\left(
\omega(a,x),
\left(
\delta(a,x),\sigma(\omega(a,x),y)
\right)
\right)
\right).
\]

\section{Standing hypotheses and levels of structure}
\label{sec:derived-poly-standing-hypotheses}

Throughout this chapter, let
\[
\mathcal E
\]
be a Grothendieck topos and let
\[
\mathcal C
=
\mathbf{Cat}(\mathcal E)_{JT,\infty}
\]
be the Joyal--Tierney localized \(\infty\)-category of internal categories.

All pullbacks, products, slices, dependent sums, dependent products, and
mapping objects below are interpreted in
\[
\mathcal C
\]
unless explicitly described as point-set or strict representative-level
constructions.

The chapter uses several layers of structure.

\begin{center}
\[
\begin{array}{c|c|c}
\text{level} & \text{structure} & \text{construction} \\
\hline
0 & \text{finite limits in }\mathcal C &
\text{derived correspondences and contextual execution} \\
1 & \text{derived arrow objects} &
\operatorname{Arr}(\mathbb A),\operatorname{Arr}(\mathbb Y),
\operatorname{Arr}(\mathbb Z) \\
2 & \text{right fibrations stable under pullback and composition} &
\text{cartesian contextual transport} \\
3 & \text{local dependent product }\Pi_{\kappa_\rho} &
\text{derived response functor }\mathcal M_\rho \\
4 & \text{regular derived predicate theory} &
\diamond\text{-modalities and existential response liftings} \\
5 & \text{first-order Heyting derived predicate theory} &
\Box\text{-modalities and profile boxes} \\
6 & \text{coherent finite joins} &
\text{finite guarded choice} \\
7 & \text{path-star and }\omega\text{-iterative hypotheses} &
\text{while programs and }\mu,\nu\text{ comparisons}
\end{array}
\]
\end{center}

The purely polynomial construction uses only the first four rows.  The modal
and fixed-point sections use the predicate-theoretic hypotheses introduced in
the previous derived dynamic proarrow-logic chapter.

\begin{remark}[Intrinsic versus objectwise arrow data]
\label{rem:derived-poly-intrinsic-objectwise-arrow-data}
The notation
\[
\operatorname{Arr}(\mathbb A)
\]
denotes a derived arrow category, equivalently a derived internal functor
category from the walking arrow.  It is not, intrinsically, merely the discrete
object of arrows of a strict internal category.

Consequently, the dependent products appearing below are homotopy-coherent
dependent products over derived query categories.  In split strict
presentations, their objectwise shadows recover the familiar fibrewise
products over ordinary arrows
\[
a:u\to v.
\]
The strict fibre formulae in this chapter should therefore be read as
representative-level objectwise specializations of the intrinsic derived
construction.
\end{remark}

\section{Derived arrows and right fibrations}
\label{sec:derived-poly-right-fibrations}

For a derived internal category
\[
\mathbb A\in\mathcal C,
\]
write
\[
\operatorname{Arr}(\mathbb A)
\]
for its derived arrow category.  It comes with source and target maps
\[
s_{\mathbb A},t_{\mathbb A}:
\operatorname{Arr}(\mathbb A)\to\mathbb A.
\]
In a strict internal-category presentation, the objectwise shadow of
\[
\operatorname{Arr}(\mathbb A)
\]
is presented by the ordinary arrow object
\[
A_1
\]
with source and target maps
\[
s_A,t_A:A_1\to A_0.
\]

\begin{definition}[Derived right fibration]
\label{def:derived-poly-right-fibration}
A morphism
\[
p:\mathbb E\to\mathbb A
\]
in
\[
\mathcal C
\]
is a \textbf{derived right fibration} if the canonical map
\[
\operatorname{Arr}(\mathbb E)
\longrightarrow
\operatorname{Arr}(\mathbb A)
\times_{t_{\mathbb A},\,\mathbb A,\,p}
\mathbb E
\]
is an equivalence in
\[
\mathcal C.
\]
Equivalently, an arrow
\[
a:u\to p(e)
\]
of
\[
\mathbb A
\]
and an object
\[
e\in\mathbb E
\]
over its target determine a cartesian lift
\[
a^\ast e\to e
\]
whose space of choices is contractible, and these lifts are
homotopy-coherently compatible with identities and composition.
\end{definition}

\begin{remark}[Strict presentations]
\label{rem:derived-poly-strict-right-fibration-presentations}
A split internal discrete fibration presents a derived right fibration.  In a
strict split presentation, the defining equivalence above is represented by
the ordinary pullback identification
\[
E_1
\cong
A_1\times_{t_A,A_0,p}E_0.
\]
Thus ordinary right actions of internal categories are strict presentations of
derived right fibrations.
\end{remark}

\begin{lemma}[Stability of derived right fibrations]
\label{lem:derived-poly-right-fibration-stability}
Derived right fibrations are stable under homotopy pullback and under
composition.
\end{lemma}

\begin{proof}
The defining condition says that the arrow object of the total category is
obtained, up to equivalence, by pulling back the arrow object of the base along
the target map.  This condition is preserved by base change.  For composites,
the cartesian lifts are obtained by composing the cartesian lifts for the two
right fibrations.  Contractibility of the lift spaces and the coherent identity
and composition laws are inherited from the two factors.
\end{proof}

\begin{definition}[Right-fibrational derived Mealy data]
\label{def:derived-poly-right-fibrational-mealy-data}
A \textbf{right-fibrational derived Mealy morphism}
\[
\mathbb A\rightsquigarrow^h\mathbb B
\]
is a derived Mealy morphism
\[
\chi=(\pi,\Omega):
\mathbb X\to\mathbb A\times\mathbb B
\]
such that
\[
\pi:\mathbb X\to\mathbb A
\]
is a derived right fibration.

A \textbf{right-fibrational downstream behaviour} over
\[
\mathbb B
\]
is a downstream behaviour
\[
\rho:\mathbb Y\to\mathbb B
\]
which is a derived right fibration.
\end{definition}

Thus the input side of a right-fibrational derived Mealy morphism carries
homotopy-coherent contravariant transport along input arrows, and the
downstream behaviour carries homotopy-coherent contravariant transport along
output arrows.

\section{Contextual right fibrations}
\label{sec:derived-poly-contextual-right-fibration}

Let
\[
\chi=(\pi,\Omega):\mathbb X\to\mathbb A\times\mathbb B
\]
be a right-fibrational derived Mealy morphism, and let
\[
\rho:\mathbb Y\to\mathbb B
\]
be a right-fibrational downstream behaviour.

Define the contextual execution object by the homotopy pullback
\[
\mathbb Z
:=
\mathbb X\times_{\mathbb B}\mathbb Y
\]
in
\[
\mathcal C.
\]
It comes with projections
\[
\operatorname{pr}_{\mathbb X}:\mathbb Z\to\mathbb X,
\qquad
\operatorname{pr}_{\mathbb Y}:\mathbb Z\to\mathbb Y.
\]
The contextual input projection is
\[
I_{\chi,\rho}:
\mathbb Z\to\mathbb A,
\qquad
I_{\chi,\rho}:=\pi\operatorname{pr}_{\mathbb X}.
\]

\begin{proposition}[Contextual execution is right-fibrational]
\label{prop:derived-poly-contextual-right-fibration}
The morphism
\[
I_{\chi,\rho}:\mathbb Z\to\mathbb A
\]
is a derived right fibration.
\end{proposition}

\begin{proof}
The map
\[
\mathbb Z\to\mathbb X
\]
is the homotopy pullback of
\[
\rho:\mathbb Y\to\mathbb B
\]
along
\[
\Omega:\mathbb X\to\mathbb B.
\]
By Lemma~\ref{lem:derived-poly-right-fibration-stability},
\[
\mathbb Z\to\mathbb X
\]
is a derived right fibration.  Since
\[
\pi:\mathbb X\to\mathbb A
\]
is a derived right fibration and derived right fibrations are closed under
composition, the composite
\[
I_{\chi,\rho}:\mathbb Z\to\mathbb A
\]
is a derived right fibration.

Equivalently, if
\[
z=(x,y)\in\mathbb Z
\]
and
\[
a:u\to I_{\chi,\rho}(z)=\pi(x)
\]
is an input arrow, then cartesian transport in
\[
\pi:\mathbb X\to\mathbb A
\]
gives, up to contractible choice, a lift
\[
\xi:x'\to x
\]
over
\[
a.
\]
The output cocycle gives an arrow
\[
\Omega(\xi):\Omega(x')\to\Omega(x)
\]
of
\[
\mathbb B.
\]
Since
\[
\Omega(x)=\rho(y),
\]
cartesian transport in
\[
\rho:\mathbb Y\to\mathbb B
\]
gives, again up to contractible choice, a lift
\[
\eta:y'\to y
\]
over
\[
\Omega(\xi).
\]
Then
\[
z'=(x',y')
\]
lies in
\[
\mathbb Z
\]
and
\[
(\xi,\eta):z'\to z
\]
is the cartesian lift of
\[
a
\]
through
\[
I_{\chi,\rho}.
\]
\end{proof}

\begin{remark}[Contextual transport]
\label{rem:derived-poly-contextual-transport}
The proof describes the one-step contextual execution operation.  A current
contextual state
\[
z=(x,y)
\]
and an input arrow
\[
a:u\to\pi(x)
\]
determine a lifted input execution
\[
\xi:x'\to x,
\]
a downstream response
\[
\eta:y'\to y,
\]
and a next contextual state
\[
a^\ast z=(x',y').
\]
The polynomial response coalgebra below is the dependent-product transpose of
this cartesian transport operation.
\end{remark}

\section{The derived interface, queries, and responses}
\label{sec:derived-poly-interface-queries-responses}

The derived interface is
\[
\mathbb J_\rho
:=
\mathbb A\times\mathbb Y.
\]
The contextual execution object
\[
\mathbb Z=\mathbb X\times_{\mathbb B}\mathbb Y
\]
maps to
\[
\mathbb J_\rho
\]
by
\[
j_{\chi,\rho}:\mathbb Z\to\mathbb A\times\mathbb Y,
\]
whose two components are
\[
I_{\chi,\rho}:\mathbb Z\to\mathbb A
\]
and
\[
\operatorname{pr}_{\mathbb Y}:\mathbb Z\to\mathbb Y.
\]

\subsection{Input queries}

Define the derived query object by
\[
\mathbb Q_\rho
:=
\operatorname{Arr}(\mathbb A)
\times_{t_{\mathbb A},\,\mathbb A,\,\operatorname{pr}_{\mathbb A}}
\mathbb J_\rho.
\]
Thus a query is a derived input arrow whose target is the current input
component of the interface, together with the current downstream object.

Write
\[
\kappa_\rho:\mathbb Q_\rho\to\mathbb J_\rho
\]
for the projection to the current interface.  There is also a projection
\[
\operatorname{in}:\mathbb Q_\rho\to\operatorname{Arr}(\mathbb A),
\]
and the next input component is
\[
s_{\mathbb A}\operatorname{in}:\mathbb Q_\rho\to\mathbb A.
\]

\subsection{Downstream responses}

A response to a query is a downstream arrow whose target is the downstream
component of the current interface.  Define the response object by
\[
\mathbb O_\rho
:=
\mathbb Q_\rho
\times_{\operatorname{pr}_{\mathbb Y}\kappa_\rho,\,\mathbb Y,\,t_{\mathbb Y}}
\operatorname{Arr}(\mathbb Y).
\]
Write
\[
r_\rho:\mathbb O_\rho\to\mathbb Q_\rho
\]
for the projection.

There are canonical label maps
\[
\operatorname{in}:\mathbb O_\rho\to\operatorname{Arr}(\mathbb A),
\]
\[
\operatorname{down}:\mathbb O_\rho\to\operatorname{Arr}(\mathbb Y),
\]
and
\[
\operatorname{out}:\mathbb O_\rho\to\operatorname{Arr}(\mathbb B),
\]
where
\[
\operatorname{out}
=
\operatorname{Arr}(\rho)\operatorname{down}.
\]

The continuation-interface map is
\[
n_\rho:\mathbb O_\rho\to\mathbb J_\rho=\mathbb A\times\mathbb Y
\]
defined by
\[
n_\rho
=
\left(
s_{\mathbb A}\operatorname{in},
s_{\mathbb Y}\operatorname{down}
\right).
\]

Thus a response records an input arrow, a downstream response arrow, the raw
output arrow induced by
\[
\rho,
\]
and the continuation interface.

\begin{remark}[Strict objectwise reading]
\label{rem:derived-poly-response-strict-reading}
In a strict split presentation of a right
\[
\mathbb B
\]
-action
\[
Y\xrightarrow{r}B_0,
\]
the derived right fibration
\[
Y\rtimes\mathbb B\to\mathbb B
\]
has arrows represented objectwise by pairs
\[
(\beta,y)
\]
where
\[
\beta:b'\to r(y),
\qquad
y'=\sigma(\beta,y).
\]
Thus an arrow
\[
\eta:y'\to y
\]
in the downstream action category is represented by the output arrow
\[
\beta
\]
and the current downstream state
\[
y.
\]
With this representative-level identification, the intrinsic derived response
object specializes objectwise to the strict response object of the
non-derived polynomial chapter.
\end{remark}

\section{The local derived response functor}
\label{sec:derived-poly-response-polynomial}

The polynomial construction below is local to the query map
\[
\kappa_\rho:\mathbb Q_\rho\to\mathbb J_\rho.
\]
It is defined whenever the corresponding dependent product exists in the
derived slice.  The dependent sum
\[
\Sigma_{r_\rho}
\]
is postcomposition along
\[
r_\rho:\mathbb O_\rho\to\mathbb Q_\rho.
\]

\begin{definition}[Derived contextual response functor]
\label{def:derived-poly-response-functor}
The \textbf{derived contextual response functor} associated to
\[
(\mathbb A,\rho)
\]
is the endofunctor
\[
\mathcal M_\rho:
\mathcal C/\mathbb J_\rho\to\mathcal C/\mathbb J_\rho
\]
defined by the local polynomial composite
\[
\mathcal M_\rho
:=
\Pi_{\kappa_\rho}\Sigma_{r_\rho}n_\rho^\ast,
\]
whenever the displayed dependent product exists.
\end{definition}

For an object
\[
X\to\mathbb J_\rho,
\]
the object
\[
\mathcal M_\rho(X)
\]
classifies homotopy-coherent total response profiles: for every input query at
the current interface, one chooses a downstream response and a continuation
state over the resulting continuation interface.

In split strict presentations, the objectwise shadow has the familiar
fibrewise form
\[
\mathcal M_\rho(X)_{(v,y)}
\cong
\prod_{a:u\to v}
\sum_{\eta:y'\to y}
X_{(u,y')}.
\]
Equivalently, writing a downstream arrow
\[
\eta:y'\to y
\]
as a pair
\[
(\beta,y)
\]
with
\[
\beta:b'\to r(y),
\qquad
y'=\sigma(\beta,y),
\]
one obtains the strict fibre formula
\[
\prod_{a:u\to v}
\sum_{\beta:b'\to r(y)}
X_{(u,\sigma(\beta,y))}.
\]

\begin{remark}[Homotopy-coherent profile semantics]
\label{rem:derived-poly-homotopy-coherent-profiles}
The displayed product--sum expression is not the intrinsic definition of
\[
\mathcal M_\rho.
\]
It is the objectwise shadow obtained from a split strict representative.  The
intrinsic functor
\[
\Pi_{\kappa_\rho}\Sigma_{r_\rho}n_\rho^\ast
\]
retains the higher morphisms among derived queries and derived responses, and
therefore should be read as a homotopy-coherent enhancement of the strict
polynomial response functor.
\end{remark}

\begin{remark}[Dependence]
\label{rem:derived-poly-response-functor-dependence}
The response functor
\[
\mathcal M_\rho
\]
depends on the input object
\[
\mathbb A
\]
and the downstream right fibration
\[
\rho:\mathbb Y\to\mathbb B.
\]
The Mealy morphism
\[
\chi
\]
enters through the canonical coalgebra
\[
c_{\chi,\rho}:\mathbb Z\to\mathcal M_\rho(\mathbb Z),
\]
constructed below.
\end{remark}

\section{The canonical response coalgebra}
\label{sec:derived-poly-canonical-response-coalgebra}

Let
\[
\chi=(\pi,\Omega):\mathbb X\to\mathbb A\times\mathbb B
\]
be a right-fibrational derived Mealy morphism, and let
\[
\rho:\mathbb Y\to\mathbb B
\]
be a right-fibrational downstream behaviour.  Let
\[
\mathbb Z:=\mathbb X\times_{\mathbb B}\mathbb Y.
\]
The map
\[
j_{\chi,\rho}:\mathbb Z\to\mathbb J_\rho
\]
makes
\[
\mathbb Z
\]
an object of
\[
\mathcal C/\mathbb J_\rho.
\]

The cartesian transport construction of
Proposition~\ref{prop:derived-poly-contextual-right-fibration} defines an
uncurried map over
\[
\mathbb Q_\rho:
\]
\[
\widetilde c_{\chi,\rho}:
\kappa_\rho^\ast\mathbb Z
\longrightarrow
\Sigma_{r_\rho}n_\rho^\ast\mathbb Z.
\]
Indeed, an object of
\[
\kappa_\rho^\ast\mathbb Z
\]
is a contextual state
\[
z=(x,y)
\]
together with an input query
\[
a:u\to I_{\chi,\rho}(z).
\]
The right fibration
\[
\pi:\mathbb X\to\mathbb A
\]
transports
\[
x
\]
along
\[
a
\]
to a cartesian arrow
\[
\xi:x'\to x.
\]
The output cocycle gives
\[
\Omega(\xi):\Omega(x')\to\Omega(x).
\]
The right fibration
\[
\rho:\mathbb Y\to\mathbb B
\]
then transports
\[
y
\]
along
\[
\Omega(\xi)
\]
to a cartesian arrow
\[
\eta:y'\to y.
\]
The result is the downstream response
\[
\eta
\]
together with the continuation contextual state
\[
z'=(x',y').
\]
Since the spaces of cartesian lifts are contractible, this construction is
canonical in the derived sense.

\begin{definition}[Canonical derived response coalgebra]
\label{def:derived-poly-canonical-response-coalgebra}
The \textbf{canonical derived response coalgebra} of
\[
(\chi,\rho)
\]
is the transpose of
\[
\widetilde c_{\chi,\rho}
\]
under
\[
\kappa_\rho^\ast\dashv\Pi_{\kappa_\rho}:
\]
\[
c_{\chi,\rho}:
\mathbb Z
\longrightarrow
\Pi_{\kappa_\rho}\Sigma_{r_\rho}n_\rho^\ast\mathbb Z
=
\mathcal M_\rho(\mathbb Z).
\]
\end{definition}

\begin{theorem}[Coalgebra as transpose of cartesian execution]
\label{thm:derived-poly-coalgebra-transpose}
The coalgebra
\[
c_{\chi,\rho}
\]
is the dependent-product transpose of contextual cartesian transport through
\[
I_{\chi,\rho}:\mathbb Z\to\mathbb A.
\]
Equivalently, it is the curried total-response form of the one-step execution
operation
\[
(z,a)
\longmapsto
(\eta,a^\ast z),
\]
where
\[
a^\ast z
\]
is obtained by first transporting through
\[
\pi:\mathbb X\to\mathbb A
\]
and then transporting through
\[
\rho:\mathbb Y\to\mathbb B
\]
along the output arrow induced by
\[
\Omega.
\]
\end{theorem}

\begin{proof}
The uncurried operation is precisely the map
\[
\widetilde c_{\chi,\rho}:
\kappa_\rho^\ast\mathbb Z
\to
\Sigma_{r_\rho}n_\rho^\ast\mathbb Z.
\]
The adjunction
\[
\kappa_\rho^\ast\dashv\Pi_{\kappa_\rho}
\]
identifies such maps with maps
\[
\mathbb Z
\to
\Pi_{\kappa_\rho}\Sigma_{r_\rho}n_\rho^\ast\mathbb Z.
\]
The transpose is, by definition,
\[
c_{\chi,\rho}.
\]
\end{proof}

\section{Multiplicativity from cartesian transport}
\label{sec:derived-poly-multiplicativity}

The strict polynomial chapter described multiplicativity by equations.  In
the derived setting, the canonical coalgebra is multiplicative because it is
built from cartesian transport in a right fibration.

\begin{proposition}[Canonical multiplicativity]
\label{prop:derived-poly-canonical-coalgebra-multiplicative}
The canonical response coalgebra
\[
c_{\chi,\rho}
\]
is homotopy-coherently compatible with identity and composite input queries.
Equivalently, after flattening, the forward execution correspondence
\[
\mathbb Z
\xleftarrow{t_{\mathbb Z}}
\operatorname{Arr}(\mathbb Z)
\xrightarrow{s_{\mathbb Z}}
\mathbb Z
\]
carries the monoid structure in
\[
\operatorname{Corr}(\mathcal C)(\mathbb Z,\mathbb Z)
\]
induced by the opposite-oriented composition of arrows in
\[
\mathbb Z.
\]
\end{proposition}

\begin{proof}
The uncurried form of
\[
c_{\chi,\rho}
\]
is cartesian transport in the derived right fibration
\[
I_{\chi,\rho}:\mathbb Z\to\mathbb A.
\]
Cartesian transport in a right fibration is compatible with identity arrows
and composition, coherently up to contractible choice.  Therefore its
dependent-product transpose is multiplicative in the derived sense.

The arrows of the right-fibrational action have backward orientation
\[
a^\ast z\to z.
\]
The dynamic execution correspondence reads these arrows in the forward
direction
\[
z\to a^\ast z.
\]
Hence the induced monoid structure on the flattened forward correspondence is
the one obtained from opposite-oriented arrow composition.
\end{proof}

In split strict presentations, this says the following.  If
\[
c(z)(b)=(\eta_b,z_b)
\]
and
\[
c(z_b)(a)=(\eta_a,z_{ab}),
\]
then
\[
c(z)(m_A(a,b))
=
(m_Y(\eta_a,\eta_b),z_{ab}).
\]
Here
\[
m_A(a,b)=b\circ a,
\]
and
\[
m_Y(\eta_a,\eta_b)=\eta_b\circ\eta_a
\]
with the composition convention fixed in the previous chapters.

\section{Position objects and flattening}
\label{sec:derived-poly-position-flattening}

The polynomial coalgebra is the curried form of one-step execution.  Flattening
recovers the execution correspondence.

Let
\[
H_{\mathbb Z}
:=
\Sigma_{r_\rho}n_\rho^\ast\mathbb Z
\]
as an object of
\[
\mathcal C/\mathbb Q_\rho.
\]
Then
\[
\mathcal M_\rho(\mathbb Z)
=
\Pi_{\kappa_\rho}H_{\mathbb Z}.
\]

\begin{definition}[Position object]
\label{def:derived-poly-position-object}
The \textbf{position object} of
\[
\mathcal M_\rho
\]
at
\[
\mathbb Z
\]
is
\[
\mathrm{Pos}_{\rho}(\mathbb Z)
:=
\kappa_\rho^\ast\Pi_{\kappa_\rho}H_{\mathbb Z}.
\]
It comes with the projection
\[
\rho_{\mathrm{pos}}:
\mathrm{Pos}_{\rho}(\mathbb Z)
\to
\mathcal M_\rho(\mathbb Z),
\]
and, by the counit of
\[
\kappa_\rho^\ast\dashv\Pi_{\kappa_\rho},
\]
an evaluation map
\[
\operatorname{ev}:
\mathrm{Pos}_{\rho}(\mathbb Z)
\to
H_{\mathbb Z}.
\]
\end{definition}

Unpacking
\[
H_{\mathbb Z}=\Sigma_{r_\rho}n_\rho^\ast\mathbb Z,
\]
the evaluation map yields canonical maps
\[
\operatorname{in}:
\mathrm{Pos}_{\rho}(\mathbb Z)\to\operatorname{Arr}(\mathbb A),
\]
\[
\operatorname{down}:
\mathrm{Pos}_{\rho}(\mathbb Z)\to\operatorname{Arr}(\mathbb Y),
\]
\[
\operatorname{out}:
\mathrm{Pos}_{\rho}(\mathbb Z)\to\operatorname{Arr}(\mathbb B),
\]
and
\[
\operatorname{next}:
\mathrm{Pos}_{\rho}(\mathbb Z)\to\mathbb Z.
\]
The output label is
\[
\operatorname{out}
=
\operatorname{Arr}(\rho)\operatorname{down}.
\]

\begin{definition}[Flattened execution object]
\label{def:derived-poly-flattened-execution-object}
The \textbf{flattened execution object} associated to
\[
c_{\chi,\rho}
\]
is the homotopy pullback
\[
\mathbb M_{\chi,\rho}
:=
\mathbb Z
\times_{\mathcal M_\rho(\mathbb Z)}
\mathrm{Pos}_{\rho}(\mathbb Z).
\]
It carries maps
\[
s_{\mathrm{for}}:\mathbb M_{\chi,\rho}\to\mathbb Z
\]
and
\[
t_{\mathrm{for}}:\mathbb M_{\chi,\rho}\to\mathbb Z,
\]
where
\[
s_{\mathrm{for}}
\]
is the projection to the current contextual state and
\[
t_{\mathrm{for}}
\]
is induced by
\[
\operatorname{next}:
\mathrm{Pos}_{\rho}(\mathbb Z)\to\mathbb Z.
\]
Thus flattening gives a forward execution correspondence
\[
\mathbf m_{\chi,\rho}
=
\left(
\mathbb Z
\xleftarrow{s_{\mathrm{for}}}
\mathbb M_{\chi,\rho}
\xrightarrow{t_{\mathrm{for}}}
\mathbb Z
\right).
\]
\end{definition}

\begin{theorem}[Flattening identifies with the arrow object]
\label{thm:derived-poly-flattening-arrow-object}
There is a canonical equivalence
\[
\mathbb M_{\chi,\rho}
\simeq
\operatorname{Arr}(\mathbb Z).
\]
Under this equivalence, the flattened forward correspondence is
\[
\mathbb Z
\xleftarrow{t_{\mathbb Z}}
\operatorname{Arr}(\mathbb Z)
\xrightarrow{s_{\mathbb Z}}
\mathbb Z.
\]
\end{theorem}

\begin{proof}
A point of
\[
\mathbb M_{\chi,\rho}
\]
is a contextual state
\[
z
\]
together with a chosen input-query position in the response profile
\[
c_{\chi,\rho}(z).
\]
Equivalently, it is an input arrow
\[
a:u\to I_{\chi,\rho}(z)
\]
together with the cartesian lift of
\[
a
\]
through
\[
I_{\chi,\rho}:\mathbb Z\to\mathbb A.
\]
Because
\[
I_{\chi,\rho}
\]
is a derived right fibration, such data are canonically equivalent to an arrow
of
\[
\mathbb Z.
\]
Thus
\[
\mathbb M_{\chi,\rho}
\simeq
\operatorname{Arr}(\mathbb Z).
\]

An arrow of the right-fibrational action is oriented contravariantly:
\[
a^\ast z\to z.
\]
The forward execution step is read as
\[
z\to a^\ast z.
\]
Therefore the forward execution correspondence has source
\[
t_{\mathbb Z}:\operatorname{Arr}(\mathbb Z)\to\mathbb Z
\]
and target
\[
s_{\mathbb Z}:\operatorname{Arr}(\mathbb Z)\to\mathbb Z.
\]
\end{proof}

\begin{remark}[Orientation]
\label{rem:derived-poly-orientation-forward-backward}
The contextual right fibration
\[
I_{\chi,\rho}:\mathbb Z\to\mathbb A
\]
is contravariant.  Its arrows have backward action orientation
\[
a^\ast z\to z.
\]
Dynamic programs are read in the forward orientation
\[
z\to a^\ast z.
\]
Thus the flattened program correspondence uses the opposite-oriented arrow
span
\[
\mathbb Z
\xleftarrow{t_{\mathbb Z}}
\operatorname{Arr}(\mathbb Z)
\xrightarrow{s_{\mathbb Z}}
\mathbb Z.
\]
This is the intrinsic form of the strict comparison
\[
(P\rtimes\mathbb A)\times_{\mathbb B}(Y\rtimes\mathbb B)
\cong
(\mathsf{Prog}_{\Phi,Y})^{\mathrm{op}}.
\]
\end{remark}

\section{Guard loci from intrinsic labels}
\label{sec:derived-poly-guard-loci}

Let
\[
\mathbf m_{\chi,\rho}:\mathbb Z\nrightarrow\mathbb Z
\]
be the flattened forward execution correspondence.  Using the equivalence
\[
\mathbb M_{\chi,\rho}
\simeq
\operatorname{Arr}(\mathbb Z),
\]
its apex carries canonical maps
\[
\iota_A:
\operatorname{Arr}(\mathbb Z)\to\operatorname{Arr}(\mathbb A),
\]
\[
\iota_Y:
\operatorname{Arr}(\mathbb Z)\to\operatorname{Arr}(\mathbb Y),
\]
and
\[
\iota_B:
\operatorname{Arr}(\mathbb Z)\to\operatorname{Arr}(\mathbb B).
\]
Here
\[
\iota_A
\]
records the input arrow,
\[
\iota_Y
\]
records the downstream response arrow, and
\[
\iota_B
=
\operatorname{Arr}(\rho)\iota_Y
\]
records the raw output arrow.

Thus guarded primitive specifications may involve predicates
\[
\theta\in\mathrm{Pred}(\mathbb Z),
\]
\[
R\in\mathrm{Pred}(\operatorname{Arr}(\mathbb A)),
\]
\[
O\in\mathrm{Pred}(\operatorname{Arr}(\mathbb B)),
\]
\[
T\in\mathrm{Pred}(\operatorname{Arr}(\mathbb Y)),
\]
and
\[
\varphi\in\mathrm{Pred}(\mathbb Z),
\]
where
\[
\mathrm{Pred}
\]
is the derived predicate theory fixed in the previous derived proarrow-logic
chapter.

\begin{remark}[Strict guard predicates]
\label{rem:derived-poly-strict-guard-predicates}
The intrinsic guards above live on derived arrow categories.  In a split strict
representative, the old non-derived guards on
\[
A_1,
\qquad
B_1,
\qquad
(\int_{\mathbb B}Y)_1
\]
are recovered by taking representative-level predicates on the objectwise
arrow parts of
\[
\operatorname{Arr}(\mathbb A),
\qquad
\operatorname{Arr}(\mathbb B),
\qquad
\operatorname{Arr}(\mathbb Y).
\]
\end{remark}

\begin{definition}[Guarded primitive contextual program]
\label{def:derived-poly-guarded-primitive}
Let
\[
\theta\in\mathrm{Pred}(\mathbb Z),
\quad
R\in\mathrm{Pred}(\operatorname{Arr}(\mathbb A)),
\quad
O\in\mathrm{Pred}(\operatorname{Arr}(\mathbb B)),
\quad
T\in\mathrm{Pred}(\operatorname{Arr}(\mathbb Y)).
\]
The guarded primitive program
\[
\theta;R/O/T
\]
is the subcorrespondence of
\[
\mathbf m_{\chi,\rho}
\]
whose apex predicate is
\[
t_{\mathbb Z}^{\ast}\theta
\wedge
\iota_A^\ast R
\wedge
\iota_B^\ast O
\wedge
\iota_Y^\ast T
\]
on
\[
\operatorname{Arr}(\mathbb Z).
\]
Equivalently,
\[
\theta;R/O/T
\]
is the restricted span
\[
\mathbb Z
\xleftarrow{t_{\mathbb Z}i}
(\theta;R/O/T)^\sharp
\xrightarrow{s_{\mathbb Z}i}
\mathbb Z,
\]
where
\[
i:(\theta;R/O/T)^\sharp\hookrightarrow\operatorname{Arr}(\mathbb Z)
\]
is the displayed subobject.
\end{definition}

\section{Guarded polynomial response liftings}
\label{sec:derived-poly-response-liftings}

We now use the regular fragment of the derived predicate theory fixed in the
previous chapter, so that dependent sums and Beck--Chevalley are available for
the displayed maps.

Let
\[
\rho_{\mathrm{pos}}:
\mathrm{Pos}_{\rho}(\mathbb Z)
\to
\mathcal M_\rho(\mathbb Z)
\]
be the position projection.

\begin{definition}[Guarded existential response lifting]
\label{def:derived-poly-guarded-existential-lifting}
For predicates
\[
R\in\mathrm{Pred}(\operatorname{Arr}(\mathbb A)),
\]
\[
O\in\mathrm{Pred}(\operatorname{Arr}(\mathbb B)),
\]
\[
T\in\mathrm{Pred}(\operatorname{Arr}(\mathbb Y)),
\]
and
\[
\varphi\in\mathrm{Pred}(\mathbb Z),
\]
define the guarded existential response lifting
\[
\lambda_{R/O/T}(\varphi)
\in
\mathrm{Pred}(\mathcal M_\rho(\mathbb Z))
\]
by
\[
\lambda_{R/O/T}(\varphi)
:=
\Sigma_{\rho_{\mathrm{pos}}}
\left(
\operatorname{in}^{\ast}R
\wedge
\operatorname{out}^{\ast}O
\wedge
\operatorname{down}^{\ast}T
\wedge
\operatorname{next}^{\ast}\varphi
\right).
\]
\end{definition}

This predicate says that a total response profile has some input position
whose input label satisfies
\[
R,
\]
whose raw output satisfies
\[
O,
\]
whose downstream response satisfies
\[
T,
\]
and whose continuation satisfies
\[
\varphi.
\]

\begin{theorem}[Polynomial lifting/span diamond comparison]
\label{thm:derived-poly-lifting-span-comparison}
For every state predicate
\[
\theta\in\mathrm{Pred}(\mathbb Z)
\]
and every postcondition
\[
\varphi\in\mathrm{Pred}(\mathbb Z),
\]
one has
\[
\theta
\wedge
c_{\chi,\rho}^{\ast}\lambda_{R/O/T}(\varphi)
=
\langle\theta;R/O/T\rangle\varphi.
\]
\end{theorem}

\begin{proof}
By definition,
\[
\lambda_{R/O/T}(\varphi)
=
\Sigma_{\rho_{\mathrm{pos}}}
\left(
\operatorname{in}^{\ast}R
\wedge
\operatorname{out}^{\ast}O
\wedge
\operatorname{down}^{\ast}T
\wedge
\operatorname{next}^{\ast}\varphi
\right).
\]
Pulling back along
\[
c_{\chi,\rho}
\]
and using Beck--Chevalley identifies the pullback of
\[
\rho_{\mathrm{pos}}
\]
with
\[
\operatorname{Arr}(\mathbb Z)
\]
by Theorem~\ref{thm:derived-poly-flattening-arrow-object}.  Under this
identification,
\[
\operatorname{in}=\iota_A,
\qquad
\operatorname{out}=\iota_B,
\qquad
\operatorname{down}=\iota_Y,
\qquad
\operatorname{next}=s_{\mathbb Z}.
\]
Therefore
\[
c_{\chi,\rho}^{\ast}\lambda_{R/O/T}(\varphi)
=
\Sigma_{t_{\mathbb Z}}
\left(
\iota_A^\ast R
\wedge
\iota_B^\ast O
\wedge
\iota_Y^\ast T
\wedge
s_{\mathbb Z}^{\ast}\varphi
\right).
\]
By Frobenius reciprocity,
\[
\theta
\wedge
c_{\chi,\rho}^{\ast}\lambda_{R/O/T}(\varphi)
=
\Sigma_{t_{\mathbb Z}}
\left(
t_{\mathbb Z}^{\ast}\theta
\wedge
\iota_A^\ast R
\wedge
\iota_B^\ast O
\wedge
\iota_Y^\ast T
\wedge
s_{\mathbb Z}^{\ast}\varphi
\right).
\]
This is exactly the diamond of the restricted correspondence
\[
\theta;R/O/T.
\]
\end{proof}

\begin{remark}[Central comparison]
\label{rem:derived-poly-central-comparison}
Theorem~\ref{thm:derived-poly-lifting-span-comparison} is the central bridge
of the chapter.  The polynomial lifting is a predicate on total response
profiles.  Pulling it back along the response coalgebra recovers the guarded
diamond modality of the flattened execution correspondence.
\end{remark}

\section{Profile boxes}
\label{sec:derived-poly-profile-boxes}

We now use the first-order Heyting fragment of the same derived predicate
theory, so that pullback functors have right adjoints satisfying
Beck--Chevalley.

\begin{definition}[Weak universal profile box]
\label{def:derived-poly-weak-profile-box}
For predicates
\[
R\in\mathrm{Pred}(\operatorname{Arr}(\mathbb A)),
\]
\[
O\in\mathrm{Pred}(\operatorname{Arr}(\mathbb B)),
\]
\[
T\in\mathrm{Pred}(\operatorname{Arr}(\mathbb Y)),
\]
and
\[
\varphi\in\mathrm{Pred}(\mathbb Z),
\]
define
\[
\Box^{\mathrm{prof}}_{R/O/T}(\varphi)
:=
\Pi_{\rho_{\mathrm{pos}}}
\left(
(
\operatorname{in}^{\ast}R
\wedge
\operatorname{out}^{\ast}O
\wedge
\operatorname{down}^{\ast}T
)
\Rightarrow
\operatorname{next}^{\ast}\varphi
\right).
\]
\end{definition}

This says that every position of the response profile satisfying the input,
output, and downstream guards has a continuation satisfying
\[
\varphi.
\]

\begin{definition}[Strong universal profile box]
\label{def:derived-poly-strong-profile-box}
The \textbf{strong universal profile box} is
\[
\Box^{\mathrm{prof,strong}}_{R/O/T}(\varphi)
:=
\Pi_{\rho_{\mathrm{pos}}}
\left(
\operatorname{in}^{\ast}R
\Rightarrow
(
\operatorname{out}^{\ast}O
\wedge
\operatorname{down}^{\ast}T
\wedge
\operatorname{next}^{\ast}\varphi
)
\right).
\]
\end{definition}

The strong profile box requires every \(R\)-input position to produce an
output satisfying
\[
O,
\]
a downstream response satisfying
\[
T,
\]
and a continuation satisfying
\[
\varphi.
\]

\begin{theorem}[Weak profile boxes and span boxes]
\label{thm:derived-poly-weak-profile-box-span-box}
One has
\[
c_{\chi,\rho}^{\ast}
\Box^{\mathrm{prof}}_{R/O/T}(\varphi)
=
[R/O/T]\varphi,
\]
where
\[
[R/O/T]\varphi
\]
denotes the box modality of the primitive span restricted by
\[
R,
\qquad
O,
\qquad
T.
\]
More generally, for a state predicate
\[
\theta\in\mathrm{Pred}(\mathbb Z),
\]
one has
\[
\theta\Rightarrow
c_{\chi,\rho}^{\ast}
\Box^{\mathrm{prof}}_{R/O/T}(\varphi)
=
[\theta;R/O/T]\varphi.
\]
\end{theorem}

\begin{proof}
By Beck--Chevalley for the right adjoints
\[
\Pi,
\]
pulling the weak profile box back along
\[
c_{\chi,\rho}
\]
gives
\[
\Pi_{t_{\mathbb Z}}
\left(
(
\iota_A^\ast R
\wedge
\iota_B^\ast O
\wedge
\iota_Y^\ast T
)
\Rightarrow
s_{\mathbb Z}^{\ast}\varphi
\right).
\]
This is exactly the box of the restricted primitive correspondence determined
by
\[
R,
\qquad
O,
\qquad
T.
\]
For the state-guarded version, use the test-prefix law
\[
[\theta?;U]\varphi
=
\theta\Rightarrow[U]\varphi.
\]
\end{proof}

\begin{remark}[Weak and strong profile boxes]
\label{rem:derived-poly-weak-strong-profile-boxes}
The weak profile box is precisely the profile-level predicate whose pullback
along the coalgebra gives the box of the restricted execution span.

The strong profile box is generally stricter.  It requires all inputs
satisfying
\[
R
\]
to produce outputs satisfying
\[
O,
\]
downstream responses satisfying
\[
T,
\]
and continuations satisfying
\[
\varphi.
\]
It is therefore genuinely a profile-level condition rather than merely the
box of a restricted span.
\end{remark}

\section{Guarded commands, choice, conditionals, and while programs}
\label{sec:derived-poly-guarded-program-algebra}

The remaining guarded program operations are inherited from the derived
dynamic proarrow logic of the previous chapter, applied to the flattened
correspondence
\[
\mathbf m_{\chi,\rho}:\mathbb Z\nrightarrow\mathbb Z.
\]

Let
\[
U:\mathbb Z\nrightarrow\mathbb Z
\]
be an endoprogram and let
\[
\theta\in\mathrm{Pred}(\mathbb Z)
\]
be a predicate.  The test span is
\[
\theta?
=
\left(
\mathbb Z\xleftarrow{\theta}\Theta\xrightarrow{\theta}\mathbb Z
\right),
\]
and the guarded command is
\[
\theta?;U.
\]

\begin{proposition}[Guard-prefix laws]
\label{prop:derived-poly-guard-prefix-laws}
For every endoprogram
\[
U:\mathbb Z\nrightarrow\mathbb Z,
\]
one has
\[
\langle\theta?;U\rangle\varphi
=
\theta\wedge\langle U\rangle\varphi.
\]
In the first-order Heyting fragment,
\[
[\theta?;U]\varphi
=
\theta\Rightarrow[U]\varphi.
\]
\end{proposition}

\begin{proof}
These are the test laws and functoriality of diamonds and boxes:
\[
\langle\theta?\rangle\psi=\theta\wedge\psi,
\qquad
[\theta?]\psi=\theta\Rightarrow\psi.
\]
\end{proof}

Assume now that the relevant predicate fibres admit coherent finite joins and
that subprograms are interpreted as predicates on the apex of a fixed ambient
correspondence.

\begin{proposition}[Finite guarded choice]
\label{prop:derived-poly-finite-guarded-choice}
Let
\[
U,V\hookrightarrow T
\]
be subprograms inside a fixed ambient endocorrespondence
\[
T:\mathbb Z\nrightarrow\mathbb Z,
\]
and let
\[
\theta,\epsilon\in\mathrm{Pred}(\mathbb Z)
\]
be state guards.  Then
\[
\left\langle
(\theta?;U)\vee(\epsilon?;V)
\right\rangle\varphi
=
(\theta\wedge\langle U\rangle\varphi)
\vee
(\epsilon\wedge\langle V\rangle\varphi).
\]
In the first-order Heyting fragment,
\[
\left[
(\theta?;U)\vee(\epsilon?;V)
\right]\varphi
=
(\theta\Rightarrow[U]\varphi)
\wedge
(\epsilon\Rightarrow[V]\varphi).
\]
\end{proposition}

\begin{proof}
The diamond statement is finite choice for subprograms inside a common ambient
correspondence, followed by the guard-prefix diamond law.  The box statement
uses finite choice for boxes inside a common ambient correspondence and then
the guard-prefix box law.
\end{proof}

\begin{definition}[Internal conditional]
\label{def:derived-poly-internal-conditional}
Let
\[
U,V\hookrightarrow T
\]
be subprograms inside a fixed ambient endocorrespondence, and let
\[
\theta,\epsilon\in\mathrm{Pred}(\mathbb Z)
\]
be complementary guards.  Define
\[
\mathbf{if}\ \theta\ \mathbf{then}\ U\ \mathbf{else}\ V
:=
(\theta?;U)\vee(\epsilon?;V).
\]
In Boolean predicate fibres, one may take
\[
\epsilon=\neg\theta.
\]
\end{definition}

Using the path-star and fixed-point hypotheses already fixed for derived
contextual dynamic proarrow logic, one obtains the corresponding while-program
laws.

\begin{definition}[Guarded while program]
\label{def:derived-poly-guarded-while}
The guarded while program with loop guard
\[
\theta,
\]
body
\[
U,
\]
and exit guard
\[
\epsilon
\]
is
\[
\mathbf{while}\ \theta\ \mathbf{do}\ U\ \mathbf{exit}\ \epsilon
:=
(\theta?;U)^\star_{\mathrm{path}};\epsilon?.
\]
\end{definition}

\begin{theorem}[Diamond law for guarded while programs]
\label{thm:derived-poly-while-diamond}
In the \(\omega\)-iterative fragment of the derived first-order Heyting
predicate theory,
\[
\left\langle
(\theta?;U)^\star_{\mathrm{path}};\epsilon?
\right\rangle\varphi
=
\mu X.\bigl(
\epsilon\wedge\varphi
\vee
\theta\wedge\langle U\rangle X
\bigr).
\]
\end{theorem}

\begin{proof}
By functoriality of diamonds,
\[
\left\langle
(\theta?;U)^\star_{\mathrm{path}};\epsilon?
\right\rangle\varphi
=
\left\langle
(\theta?;U)^\star_{\mathrm{path}}
\right\rangle
(\epsilon\wedge\varphi).
\]
The path/fixed-point comparison gives
\[
\left\langle
(\theta?;U)^\star_{\mathrm{path}}
\right\rangle
(\epsilon\wedge\varphi)
=
\mu X.\bigl(
\epsilon\wedge\varphi
\vee
\langle\theta?;U\rangle X
\bigr).
\]
Using
\[
\langle\theta?;U\rangle X
=
\theta\wedge\langle U\rangle X
\]
gives the result.
\end{proof}

\begin{theorem}[Box law for guarded while programs]
\label{thm:derived-poly-while-box}
Under the corresponding box hypotheses,
\[
\left[
(\theta?;U)^\star_{\mathrm{path}};\epsilon?
\right]\varphi
=
\nu X.\bigl(
(\epsilon\Rightarrow\varphi)
\wedge
(\theta\Rightarrow[U]X)
\bigr).
\]
\end{theorem}

\begin{proof}
This is dual to the diamond proof, using the box test law
\[
[\epsilon?]\varphi=\epsilon\Rightarrow\varphi
\]
and the path/fixed-point comparison for boxes.
\end{proof}

\section{Domain, antidomain, and guarded automata}
\label{sec:derived-poly-domain-automata}

Let
\[
U:\mathbb Z\nrightarrow\mathbb Z
\]
be an endoprogram.  Its domain predicate is
\[
\operatorname{dom}(U)
:=
\langle U\rangle\top.
\]

\begin{proposition}[Domain-prefix law]
\label{prop:derived-poly-domain-prefix}
For every endoprogram
\[
U,
\]
one has
\[
\langle\operatorname{dom}(U)?;U\rangle\varphi
=
\langle U\rangle\varphi.
\]
\end{proposition}

\begin{proof}
By the guard-prefix law,
\[
\langle\operatorname{dom}(U)?;U\rangle\varphi
=
\operatorname{dom}(U)\wedge\langle U\rangle\varphi.
\]
Since
\[
\langle U\rangle\varphi\le\langle U\rangle\top
=
\operatorname{dom}(U),
\]
the meet is
\[
\langle U\rangle\varphi.
\]
\end{proof}

If predicate fibres are Boolean, define
\[
\operatorname{adom}(U)
:=
\neg\operatorname{dom}(U).
\]

\begin{proposition}[Antidomain annihilation]
\label{prop:derived-poly-antidomain-annihilation}
In Boolean predicate fibres,
\[
\langle\operatorname{adom}(U)?;U\rangle\varphi
=
\bot.
\]
\end{proposition}

\begin{proof}
By the guard-prefix law,
\[
\langle\operatorname{adom}(U)?;U\rangle\varphi
=
\operatorname{adom}(U)\wedge\langle U\rangle\varphi.
\]
Since
\[
\langle U\rangle\varphi\le\operatorname{dom}(U)
\]
and
\[
\operatorname{adom}(U)=\neg\operatorname{dom}(U),
\]
the meet is bottom.
\end{proof}

\begin{definition}[Derived guarded Mealy automaton]
\label{def:derived-poly-guarded-mealy-automaton}
A \textbf{derived guarded Mealy automaton} over
\[
(\chi,\rho)
\]
is a finite family of guarded primitive subprograms
\[
G_i
=
(\theta_i;R_i/O_i/T_i)^\sharp
\hookrightarrow
\operatorname{Arr}(\mathbb Z)
\]
inside the common ambient flattened execution correspondence
\[
\mathbf m_{\chi,\rho}.
\]
Its one-step program is the finite join
\[
G:=\bigvee_{i\in I}G_i.
\]
\end{definition}

\begin{proposition}[Diamond semantics of derived guarded Mealy automata]
\label{prop:derived-poly-guarded-automaton-diamond}
For a derived guarded Mealy automaton with one-step program
\[
G=\bigvee_{i\in I}G_i,
\]
one has
\[
\langle G\rangle\varphi
=
\bigvee_{i\in I}\langle G_i\rangle\varphi.
\]
\end{proposition}

\begin{proof}
This is the finite choice law for subprograms inside the common ambient
execution correspondence.
\end{proof}

A derived guarded Mealy automaton is \textbf{witness-disjoint} if
\[
G_i\wedge G_j=\bot
\qquad
(i\ne j).
\]
It is \textbf{branch-disjoint} if
\[
\operatorname{dom}(G_i)\wedge\operatorname{dom}(G_j)=\bot
\qquad
(i\ne j).
\]
It is \textbf{guard-total} if
\[
\bigvee_{i\in I}\operatorname{dom}(G_i)=\top.
\]
When branch-disjointness and guard-totality hold, the branch domains form a
finite disjoint cover of
\[
\mathbb Z.
\]
Thus the enabled branch is unique at the level of branch domains, although a
branch may still contain multiple execution witnesses.

\section{Terminal and stateless specializations}
\label{sec:derived-poly-terminal-stateless}

\subsection{Terminal trace}

Let
\[
\rho=1_{\mathbb B}:\mathbb B\to\mathbb B
\]
be the terminal downstream behaviour.  Since the identity map is a derived
right fibration, the preceding construction applies.

Then
\[
\mathbb Z
=
\mathbb X\times_{\mathbb B}\mathbb B
\simeq
\mathbb X.
\]
In a split strict presentation, the objectwise shadow of the derived response
functor has fibre
\[
\mathcal M_{1_{\mathbb B}}(X)_{(v,b)}
\cong
\prod_{a:u\to v}
\sum_{\beta:b'\to b}
X_{(u,b')}.
\]

For a strict Mealy presentation
\[
(P,\delta,\omega):\mathbb A\rightsquigarrow\mathbb B,
\]
this gives the terminal trace coalgebra
\[
c_{\Phi,\mathbf 1_{\mathbb B}}(x)(a)
=
(\omega(a,x),\delta(a,x)).
\]

\subsection{Stateless case}

Let
\[
F:\mathbb A\to\mathbb B
\]
be a derived internal functor.  It defines a stateless derived Mealy morphism
\[
\chi_F:\mathbb A\to\mathbb A\times\mathbb B
\]
classified by
\[
1_{\mathbb A}:\mathbb A\to\mathbb A
\]
and
\[
F:\mathbb A\to\mathbb B.
\]

For a right-fibrational downstream behaviour
\[
\rho:\mathbb Y\to\mathbb B,
\]
the contextual execution object is
\[
\mathbb Z
=
\mathbb A\times_{\mathbb B}\mathbb Y.
\]
The associated response coalgebra is the dependent-product transpose of
cartesian restriction of the downstream right fibration along
\[
F.
\]

In strict presentations, if
\[
Y\to B_0
\]
is a right
\[
\mathbb B
\]
-action, then
\[
c_{F,Y}(v,y)(a)
=
\left(
F_1(a),
\left(
s_A(a),
\sigma(F_1(a),y)
\right)
\right)
\]
for
\[
a:u\to v.
\]
For
\[
\rho=1_{\mathbb B},
\]
this reduces to
\[
c_F(v)(a)=(F_1(a),s_A(a)).
\]
Under the forward-opposite convention, the recovered program category is
\[
\mathbb A^{\mathrm{op}},
\]
with output labelling
\[
F^{\mathrm{op}}:\mathbb A^{\mathrm{op}}\to\mathbb B^{\mathrm{op}}.
\]

\section{Strict split presentations}
\label{sec:derived-poly-strict-presentations}

We now record the comparison with the non-derived polynomial chapter.

Let
\[
\Phi=(P,\delta,\omega):\mathbb A\rightsquigarrow\mathbb B
\]
be a strict Mealy morphism, and let
\[
Y=(Y\xrightarrow{r}B_0,\sigma)
\]
be a strict right
\[
\mathbb B
\]
-action.

The strict Mealy morphism presents the right-fibrational derived Mealy
morphism
\[
P\rtimes\mathbb A\to\mathbb A\times\mathbb B,
\]
and the strict downstream action presents the derived right fibration
\[
Y\rtimes\mathbb B\to\mathbb B.
\]

The derived contextual execution object is presented by
\[
(P\rtimes\mathbb A)\times_{\mathbb B}(Y\rtimes\mathbb B),
\]
whose object part is
\[
S_{\Phi,Y}
=
P\times_{q,B_0,r}Y.
\]
The object-level interface is
\[
A_0\times Y.
\]
The object-level query object is
\[
A_1\times_{t_A,A_0,\pi_A}(A_0\times Y),
\]
and the object-level response object is presented by
\[
Q_Y\times_{r\pi_Y\kappa_Y,B_0,t_B}B_1.
\]

An arrow
\[
\eta:y'\to y
\]
of
\[
Y\rtimes\mathbb B
\]
is represented by a pair
\[
(\beta,y)
\]
with
\[
\beta:b'\to r(y),
\qquad
y'=\sigma(\beta,y).
\]
Consequently the objectwise shadow of the intrinsic response functor has the
strict fibre
\[
\prod_{a:u\to v}
\sum_{\beta:b'\to r(y)}
X_{(u,\sigma(\beta,y))}.
\]

The canonical response coalgebra becomes
\[
c_{\Phi,Y}(x,y)(a)
=
\left(
\omega(a,x),
\left(
\delta(a,x),
\sigma(\omega(a,x),y)
\right)
\right).
\]

Flattening identifies with the arrow object of
\[
(P\rtimes\mathbb A)\times_{\mathbb B}(Y\rtimes\mathbb B).
\]
After applying the forward-opposite convention, its primitive execution span is
\[
S_{\Phi,Y}
\xleftarrow{s_{\Phi,Y}}
M_{\Phi,Y}
\xrightarrow{t_{\Phi,Y}}
S_{\Phi,Y},
\]
where
\[
M_{\Phi,Y}
=
A_1\times_{t_A,A_0,p\pi_P}S_{\Phi,Y},
\]
\[
s_{\Phi,Y}(a,x,y)=(x,y),
\]
and
\[
t_{\Phi,Y}(a,x,y)
=
\left(
\delta(a,x),
\sigma(\omega(a,x),y)
\right).
\]

The guarded lifting comparison becomes
\[
\theta\wedge
c_{\Phi,Y}^{\ast}\lambda_{R/O/T}(\varphi)
=
\langle\theta;R/O/T\rangle\varphi,
\]
and the weak profile box comparison becomes
\[
c_{\Phi,Y}^{\ast}
\Box^{\mathrm{prof}}_{R/O/T}(\varphi)
=
[R/O/T]\varphi.
\]

Thus, after choosing split strict representatives and passing to the
objectwise shadows of the derived arrow and response objects, the non-derived
polynomial chapter is recovered from the intrinsic homotopy-coherent
right-fibrational construction.

\section{Main consolidation}
\label{sec:derived-poly-main-consolidation}

\begin{theorem}[Derived polynomial transposes of contextual Mealy execution]
\label{thm:derived-poly-main-consolidation}
Let
\[
\mathcal C
=
\mathbf{Cat}(\mathcal E)_{JT,\infty}.
\]
Let
\[
\chi=(\pi,\Omega):\mathbb X\to\mathbb A\times\mathbb B
\]
be a right-fibrational derived Mealy morphism, and let
\[
\rho:\mathbb Y\to\mathbb B
\]
be a right-fibrational downstream behaviour.  Suppose that the derived arrow
objects used above exist, that derived right fibrations are stable under
homotopy pullback and composition, and that the local dependent product
\[
\Pi_{\kappa_\rho}
\]
needed to form
\[
\mathcal M_\rho
=
\Pi_{\kappa_\rho}\Sigma_{r_\rho}n_\rho^\ast
\]
exists.

Then:

\begin{enumerate}
\item The contextual execution object is
\[
\mathbb Z=\mathbb X\times_{\mathbb B}\mathbb Y.
\]

\item The contextual projection
\[
I_{\chi,\rho}:\mathbb Z\to\mathbb A
\]
is a derived right fibration.

\item The derived interface is
\[
\mathbb J_\rho=\mathbb A\times\mathbb Y.
\]

\item The query object is
\[
\mathbb Q_\rho
=
\operatorname{Arr}(\mathbb A)
\times_{t_{\mathbb A},\,\mathbb A,\,\operatorname{pr}_{\mathbb A}}
\mathbb J_\rho.
\]

\item The response object is
\[
\mathbb O_\rho
=
\mathbb Q_\rho
\times_{\operatorname{pr}_{\mathbb Y}\kappa_\rho,\,\mathbb Y,\,t_{\mathbb Y}}
\operatorname{Arr}(\mathbb Y).
\]

\item The raw output label is induced by
\[
\operatorname{Arr}(\rho):
\operatorname{Arr}(\mathbb Y)\to\operatorname{Arr}(\mathbb B).
\]

\item The derived response functor is
\[
\mathcal M_\rho
=
\Pi_{\kappa_\rho}\Sigma_{r_\rho}n_\rho^\ast.
\]

\item Contextual cartesian transport transposes to the canonical response
coalgebra
\[
c_{\chi,\rho}:\mathbb Z\to\mathcal M_\rho(\mathbb Z).
\]

\item This coalgebra is multiplicative because it is induced by functorial
cartesian transport.

\item Pulling the position object of
\[
\mathcal M_\rho
\]
back along
\[
c_{\chi,\rho}
\]
identifies canonically with
\[
\operatorname{Arr}(\mathbb Z).
\]

\item The flattened forward execution correspondence is
\[
\mathbb Z
\xleftarrow{t_{\mathbb Z}}
\operatorname{Arr}(\mathbb Z)
\xrightarrow{s_{\mathbb Z}}
\mathbb Z.
\]
\end{enumerate}

If, in addition, the derived predicate theory has the regular fragment used
for diamonds, then guarded existential response liftings pull back along
\[
c_{\chi,\rho}
\]
to guarded span diamonds:
\[
\theta
\wedge
c_{\chi,\rho}^{\ast}\lambda_{R/O/T}(\varphi)
=
\langle\theta;R/O/T\rangle\varphi.
\]

If the predicate theory has the first-order Heyting fragment used for boxes,
then weak profile boxes pull back along
\[
c_{\chi,\rho}
\]
to guarded span boxes:
\[
c_{\chi,\rho}^{\ast}
\Box^{\mathrm{prof}}_{R/O/T}(\varphi)
=
[R/O/T]\varphi.
\]

If the coherent finite-join, path-star, and \(\omega\)-iterative hypotheses of
the previous derived dynamic proarrow-logic chapter are also imposed, then
tests, guarded commands, finite guarded choice, conditionals, while programs,
domain operations, and guarded automata are inherited from the derived dynamic
proarrow logic of the flattened correspondence.

Finally, the terminal downstream behaviour recovers the derived terminal trace
polynomial, stateless derived Mealy morphisms classified by derived functors
\[
F:\mathbb A\to\mathbb B
\]
recover the polynomial transpose of downstream cartesian restriction along
\[
F,
\]
and split strict presentations recover the non-derived polynomial chapter at
the objectwise representative level.
\end{theorem}

\begin{proof}
Parts (1) and (2) are the definition of contextual execution and
Proposition~\ref{prop:derived-poly-contextual-right-fibration}.  Parts
(3)--(6) are the intrinsic interface, query, response, and label constructions
of Section~\ref{sec:derived-poly-interface-queries-responses}.  Part (7) is
Definition~\ref{def:derived-poly-response-functor}.  Part (8) is
Definition~\ref{def:derived-poly-canonical-response-coalgebra} and
Theorem~\ref{thm:derived-poly-coalgebra-transpose}.  Part (9) is
Proposition~\ref{prop:derived-poly-canonical-coalgebra-multiplicative}.  Parts
(10) and (11) are Theorem~\ref{thm:derived-poly-flattening-arrow-object}.

The diamond comparison is Theorem~\ref{thm:derived-poly-lifting-span-comparison}.
The box comparison is Theorem~\ref{thm:derived-poly-weak-profile-box-span-box}.
The guarded program operations are the constructions of
Sections~\ref{sec:derived-poly-guarded-program-algebra}--\ref{sec:derived-poly-domain-automata}.
The terminal and stateless claims are
Section~\ref{sec:derived-poly-terminal-stateless}.  The split strict
comparison is Section~\ref{sec:derived-poly-strict-presentations}.
\end{proof}

\begin{remark}[Final conceptual summary]
\label{rem:derived-poly-final-summary}
The derived polynomial calculus is the intrinsic curried form of
right-fibrational contextual Mealy execution.

The hierarchy is
\[
\text{right-fibrational derived Mealy data}
\quad\leadsto\quad
\text{contextual right fibration}
\quad\leadsto\quad
\text{cartesian transport}
\quad\leadsto\quad
\text{polynomial response coalgebra}
\quad\leadsto\quad
\text{flattened proarrow program}.
\]

The polynomial coalgebra records homotopy-coherent total response profiles.
Flattening it recovers the forward execution correspondence.  Guarded
polynomial liftings are profile-level predicates whose pullbacks along the
coalgebra are precisely the guarded dynamic proarrow modalities of the
flattened program semantics.

Thus the same homotopy-coherent execution semantics appears in two
presentations:
\[
\boxed{
\text{polynomial semantics}
=
\text{curried total-response semantics}
}
\]
and
\[
\boxed{
\text{proarrow semantics}
=
\text{flattened execution semantics}.
}
\]
\end{remark}